\chapter*{Preface}

\noindent I wrote this textbook primarily to provide students of Multivariate Calculus (also known as Vector Calculus or Calculus III) with the option of having a concise reference to assist them or a springboard into deeper mathematical topics for the sake of interest. When I was a freshman engineering major in college, I reached this crossroads and, in a way, chose both paths—I added a mathematics major, thereby embracing both perspectives. These notes, compiled into textbook form and titled ``Theories of Generalized Change,'' are designed to bridge theory with application, emphasizing problem-solving techniques and the essential skills needed in applied fields. \\

\noindent This course, like most others, requires a solid foundation in Calculus I and Calculus II, including an understanding of limits, derivatives, integrals, and series. Reviewing the concepts from Calculus I will be particularly helpful before diving in, and knowledge from Calculus II will provide additional clarity and depth. While familiarity with introductory multivariable concepts is beneficial, we will introduce and review key topics as needed to ensure all students have the necessary tools to succeed. This course will feel like we are revisiting basic Differential Calculus, but with a new spirit—one of generalization. We will constantly explore how to further generalize and extend our understanding of Calculus to higher dimensions. \\

\noindent The course topics cover vectors, vector-valued functions, partial derivatives, multiple integrals, vector calculus theorems, introductory topics in Partial Differential Equations (PDEs), and a bit of Complex Variables, each presented with examples relevant to engineering, physics, and other applied sciences. PDEs and, less frequently, Complex Variables are seldom discussed in such a course. However, they are intrinsically tied together and embody the spirit of generalizing calculus; not to mention, they are quite useful in their own right. Through these applications, students will develop an appreciation for the power of calculus in modeling and solving practical problems, as well as gain a taste of higher mathematics. \\

\noindent Each chapter is divided into lectures, roughly structured to fit a 90- to 120-minute session. This format allows flexibility for instructors to adapt the material to the pace and needs of the class. The sections within each lecture focus on specific concepts, promoting a gradual mastery of each topic before progressing to more advanced material. This structure also provides a convenient framework for creating assignments and assessments. In my opinion, Chapters 1–5 are absolutely crucial for a Multivariate Calculus class, while Chapters 6 and 7 are additional topics to include if there is extra time. \\

\noindent One of the most challenging aspects for students of mathematics is handling notation and the over-assumptions made by some texts. Extra care is taken to ensure that notation is clearly defined and labeled—mathematics is like learning a new language. I also make sure to write out problems in full and use examples that cover a broad range of use cases, allowing students to more easily make connections to further topics. \\

\noindent My goal as an educator with this text is to empower students to apply multivariable calculus concepts beyond the classroom, equipping them with a robust mathematical foundation for their future studies and careers. I hope this course serves as a valuable resource for both students and instructors, providing the clarity, depth, and insight needed to master Multivariate Calculus...I also crack jokes often. Enjoy!
\\
\\
\\
\\
\\
Joseph J. Babčanec
\\
\\
Adjunct Professor of Mathematics, Benedict College
\\
Lead Data Scientist, Anika Systems Inc.
\\
\\
\href{mailto:joseph.babcanec@benedict.edu}{joseph.babcanec@benedict.edu}.